% =====================================================================
%  ABSTRACT
%  Linea guida relatore: l'abstract deve contenere (a) un breve riassunto
%  della problematica, (b) una breve descrizione di cio' che e' stato fatto,
%  (c) i risultati ottenuti.
% =====================================================================

L'obesit\`a infantile rappresenta uno dei principali problemi di salute
pubblica nei paesi industrializzati e la sua prevenzione richiede che una
pluralit\`a di attori --- genitori, insegnanti e professionisti sanitari ---
possa accedere a informazioni affidabili, aggiornate e adeguate al proprio
ruolo. I modelli linguistici di grandi dimensioni (LLM) offrono un canale di
accesso conversazionale a questa conoscenza, ma in ambito sanitario il loro
impiego diretto comporta rischi rilevanti: la generazione di informazioni non
fondate sulle fonti (\emph{allucinazioni}) e la produzione di contenuti
clinicamente inappropriati, come diagnosi o indicazioni terapeutiche.

In questo lavoro \`e stato progettato e realizzato un agente conversazionale
in lingua italiana basato sul paradigma \emph{Retrieval-Augmented Generation}
(RAG) per il supporto informativo sulla prevenzione dell'obesit\`a infantile
nella fascia 0--10 anni. Il sistema \`e alimentato da un corpus di documenti
scientifici e linee guida italiane, adatta le risposte a cinque profili di
interlocutore e integra la pipeline RAG in un'architettura di sicurezza
(\emph{guardrail}) articolata su controlli di input, controlli di output e
una verifica di ancoraggio alle fonti a tempo di esecuzione. Il recupero delle
informazioni \`e affidato a un retriever ibrido che combina ricerca vettoriale
e ricerca per parole chiave.

La valutazione sperimentale, condotta con la metodologia RAGAS su un dataset
costruito ad hoc, ha riportato un valore di fedelt\`a alle fonti
(\emph{faithfulness}) pari a $0{,}934$. La robustezza del livello di sicurezza
\`e stata misurata tramite una campagna di \emph{red-teaming}: i controlli
deterministici bloccano il $78{,}8\%$ degli attacchi, mentre l'architettura a
difesa in profondit\`a neutralizza la totalit\`a degli attacchi che superano
il primo livello, senza produrre alcun output clinicamente inappropriato sui
casi considerati. La capacit\`a di adattamento al profilo, valutata mediante
classificazione cieca e una rubrica di appropriatezza, ha ottenuto un punteggio
medio di $4{,}15$ su $5$.
